% Part: many-valued-logic
% Chapter: syntax-and-semantics
% Section: formulas

\documentclass[../../../include/open-logic-section]{subfiles}

\begin{document}

\olfileid{mvl}{syn}{fml}

\olsection{\usetoken{P}{formula}}

\begin{defn}[الصيغة]
\ollabel{defn:formulas}
تعرّف مجموعة~$\Frm[L]$ من \emph{!!{formula}s} لغة القضايا~$\Lang L$
استقرائيًا كما يأتي:
\begin{enumerate}
\item كل~!!{propositional variable}~$\Obj p_i$ هو~!!{formula} ذرية.
\item كل رابط ذي~$0$ من المواضع (أي ثابت قضوي) في~$\Lang L$ هو
!!{formula} ذرية.
\item إذا كان~$\star$ رابطًا ذا~$n$ من المواضع في~$\Lang L$، حيث
$n>0$، وكانت~$!A_1$، \dots،~$!A_n$~!!{formula}s، فإن
$\star(!A_1, \dots, !A_n)$ تكون~!!a{formula}.
\tagitem{limitClause}{ولا شيء غير ذلك يكون~!!a{formula}.}{}
\end{enumerate}
إذا كان~$\star$ ذا موضع واحد، فغالبًا ما نكتب~$\star(!A_1)$ في
الصورة الأبسط~$\star !A_1$. وإذا كان~$\star$ ذا موضعين، فغالبًا ما
نكتب~$\star(!A_1,!A_2)$ في الصورة~$(!A_1 \star !A_2)$.
\end{defn}

وكالمعتاد، سنحذف غالبًا زوج الأقواس الخارجي حذفًا ضمنيًا.

\begin{ex}
  في اللغة القياسية~$\Lang{L_0}$، تكون
  $\Obj p_1 \lif (\Obj p_1 \land \lnot \Obj p_2)$ صيغة. وفي لغة
  منطق الضرب تكتب بدلًا من ذلك
  $\Obj p_1 \lif (\Obj p_1 \odot \lnot \Obj p_2)$. وإذا أضفنا
  $\triangle$ ذا الموضع الواحد إلى اللغة، حصلنا أيضًا على صيغ مثل
  $\triangle (\Obj p_1 \land \Obj p_2) \lif
  (\triangle \Obj p_1 \land \triangle \Obj p_2)$.
\end{ex}

\end{document}

